\section{泛函积分与拓扑扇区}

在格点规范理论中，规范场的空间是连通的，因此拓扑扇区的概念先验地没有良定义的含
义。然而，人们也早已知道，任何经典的连续规范场（包括多瞬子组态）都可以被格点场
任意好地逼近。于是各扇区被包含在场空间之中，但彼此并不分隔开。

利用威尔逊流，现在可以理解当晶格间距趋于零时各扇区究竟是如何被划分的。这样，
拓扑扇区的存在性原来是理论的一个动力学性质，而不是场的某种被假定或强加的连续性
的结果。


\subsection{场变换}

在有穷格点上（无论多大），变换
\begin{equation}
  U\to V=V_{t_0}
  \label{eq:fieldtrafo}
\end{equation}
是可逆的，而且实际上是场空间的一个微分同胚~\cite{TrivMaps}。因此可以在 QCD 泛函
积分中把积分变量从基本场变换到流时间 $t_0$ 处的场。如文献~\cite{TrivMaps} 中已指
出的，相应的雅可比行列式可以解析求出，并能用沿流的 Wilson 作用量表出。于是任意观
测量 $\Obs(V)$ 的期望值为
\begin{align}
  \langle\Obs\rangle&={1\over{\cal Z}}\int
  \rmD[V]\,\Obs(V)\,\rme^{-\tilde{S}(V)},
  \label{eq:fint}\\
  \tilde{S}(V)&=S(U)+{16g_0^2\over3a^2}\int_0^{t_0}\rmd t\,\Sw(V_t),
  \qquad
  \rmD[V]=\prod_{x,\mu}\rmd V(x,\mu),
  \label{eq:tSaction}
\end{align}
其中 $S(U)$ 表示变换前理论的总作用量（包括夸克行列式）。

显然，泛函积分可以这样改写为对任意流时间 $t$ 处规范场的积分。不过，把 $t$
取为参考时间 $t_0$ 是一个有趣而自然的选择。特别地，此时支配变换后积分的场便具有
理论基本低能标度量级的特征波长。


\subsection{主导场的光滑性}

给定格点规范场 $V$ 的光滑程度的定量度量是
\begin{equation}
  h=\max_{p}s_p,\qquad s_p=\Re\tr\{1-V(p)\},
  \label{eq:hmax}
\end{equation}
最大值取遍所有方格 $p$（如前，$V(p)$ 表示绕 $p$ 一周链变量的乘积）。在泛函
积分~\eqref{eq:fint} 中，$h$ 的所有可能取值都会出现；但鉴于期望值 $\langle s_p\rangle$
按 $a^4$ 比例标度，大的取值预期是不太可能的。

\begin{figure}[ht]
\centering
\includegraphics[width=0.6\textwidth]{images/prob}
\caption{给定方格 $p$ 上 $s_p$ [式~\eqref{eq:hmax}] 超过某指定值 $s$
的概率图。从上到下，三条曲线分别对应晶格间距 $a=0.1$、$0.07$ 和 $0.05$ fm。计算
在 $\SUthree$ 规范理论中用 3.1 小节描述的场系综进行。}
\label{fig:prob}
\end{figure}

由于作用量~\eqref{eq:tSaction} 中所有各项同号且偏向光滑组态，大方格值 $s_p$
被强烈压低这一点自然是合理的。事实上，当晶格间距减小时，给定方格 $p$ 上 $s_p$
大于某指定值 $s$ 的概率大致按 $a^{10}$ 下降（见图~\ref{fig:prob}）。因此，在固定物
理尺寸的有穷体积中，$h$ 值超过给定阈值的场组态的统计权重在连续极限下迅速趋于零，
而支配变换后泛函积分的场则在空间中变得一致地光滑。


\subsection{拓扑扇区的动力学分离}

很多年以前人们已经证明，若只允许满足某个光滑性条件的场，格点规范场的空间便会分
割成不连通的扇区~\cite{TopCharge,PhillipsStone}。该证明基于对拓扑荷的一个局域格
点表达式的几何构造，这个表达式取整数值并具有正确的经典连续极限。

在 QCD 的情形，满足
\begin{equation}
   h<0.067
   \label{eq:hbound}
\end{equation}
的场\footnote{Phillips 与 Stone~\cite{PhillipsStone} 证明的定理假定了一个单纯形
格点。为了能在当前情形应用它们，必须把格点的超立方体胞分割成单纯形。此时还需要
把规范场插值到新增的链上；利用文献~\cite{TopCharge,PhillipsStone}
中其他地方所用的同样插值方法即可容易地做到。}包含在几何构造所覆盖的子空间内，因
而非常像连续理论中的连续规范场那样落入各个拓扑扇区。因此，“扇区之间”的
场空间可以由不等式 $h\geq 0.067$ 来刻画。

回顾 4.2 小节关于支配泛函积分~\eqref{eq:fint} 的场之光滑性的讨论便可看出：当晶格
间距趋于零时，$h$ 值大的场在积分中的权重迅速下降，拓扑扇区正是因此浮现出来。例如，
在第 3 节报告的模拟中，满足界限~\eqref{eq:hbound} 的代表性场所占的比例从 $a=0.1$
fm 处的 $0\%$ 增至 $a=0.07$ 和 $0.05$ fm 处的 $8\%$ 和 $70\%$。由这些数据无法可靠
地确定扇区之间的场之权重的渐近行为，但图~\ref{fig:prob} 所示曲线提示：在给定物理
尺寸的格子上找到这类组态的概率按 $a^6$ 趋于零。


\subsection{拓扑磁化率}

一旦泛函积分~\eqref{eq:fint} 分裂为对实际上互不连通的各扇区的积分之和，在代表性
场系综中给场组态指派拓扑荷 $Q$ 就变得没有歧义了。前面提到的几何构造可以用来计算
拓扑荷；但鉴于第 3 节末讨论的普适性，用场张量 $G_{\mu\nu}$ 的对称格点表达式对拓
扑密度作直接离散化，预期在连续极限下会给出相同的矩 $\langle Q^n\rangle$。

在表~\ref{tab:params} 所列的格子上，按这些路线计算的二次矩 $\langle Q^2\rangle$
约为 $50$，而拓扑磁化率 $\chi_t$ 在统计误差内与晶格间距无关。常数拟合给出结果
\begin{equation}
   \chi_t^{1/4}=187.4(3.9)\,\MeV,
   \label{eq:chit}
\end{equation}
它与流时间 $t=0$ 处用手征格点 Dirac 算符及指标定理~\cite{IndThm}
得到的值 $194.5(2.4)$ MeV~\cite{ChiGW} 相差不到两个标准差。此外，正如一篇即将发表
的文章所示，这里算得的二次矩四次根与通过文献~\cite{ChiInd}
提出的普适公式的一个变体~\cite{ChiProj} 算得者之比，外推到连续极限时在 $3\%$
的统计不确定度内等于 $1$。

所有这些经验结果都支持如下猜想：泛函积分~\eqref{eq:fint} 中荷分布的矩与由指标定
理~\cite{IndThm,ChiGW,ChiInd,ChiProj} 得到的、从而也出现在手征 Ward
恒等式~\cite{ChiWI,ChiWII} 中的那些相重合。然而，尽管高度可信，这一等同仍有待理论
上的确立。


